% !TeX root = main.tex
\section{传感器系统}

\subsection{机器人传感器}

\subsubsection{传感器的定义、任务与分类}

工业机器人本身按照预先编好的程序动作，理论上并不需要感知外界。但现实生产中，很多情况让机器人必须依靠传感器才能正常工作。典型情形有：
\begin{itemize}
    \item 工件有时\textbf{无法有序摆放}，比如批量很小、型号种类多、工件本身很大，或者工件形状复杂，这些都让机器人无法按固定坐标去抓取。
    \item \textbf{工件和工件夹具的公差过大}，比如锻造和炉子作业中工件会因受热膨胀，铸件和压铸件边缘还常有毛刺，这些偏差都需要传感器去补偿。
    \item 工件\textbf{被连续移动的传送装置搬运}的场景，比如悬挂输送线或传送带，这时机器人需要具备传送带跟踪的能力，让机械臂的运动和不断移动的工件保持同步，喷漆、表面处理等工序中工件往往还会摆动或者悬挂得不太精确，进一步增加了难度。
    \item 工件的\textbf{转向角未知}，比如装配一个旋转部件时，机器人得先识别出它当前处于什么角度。
    \item \textbf{工件本身柔软易变形}，比如电缆、软管、电线、纺织品这类东西，形状不固定，抓取和装配都更依赖实时感知。
    \item 机械式的硬自动化容易出故障，用传感器可以及时发现并灵活应对这些异常，从而减少人工质检和人工干预的次数
    \item \textbf{加工精度要求超过了机器人本身能达到的绝对精度}，比如精密滚动轴承的装配、飞机侧翼的钻孔、大型机器人给飞机做清洗，这些场合机器人自身的定位精度不够用，必须靠传感器来补足。
    \item \textbf{机器人编程本身很繁琐}，这时可以借助传感器自动生成运动轨迹，或者先粗略编好轨迹，再叠加传感器做精细修正。
\end{itemize}

\textbf{传感器是一种测量值拾取器，它获取物体的属性、状态和过程信息，并把最重要的输入信号转换成合适的输出信号。}人类的感知能力常被当作衡量机器人传感能力的标杆，在一些专门场合，机器人传感器的表现甚至已经超过了人类。

传感器在机器人上除了做物理量到电信号的转换之外，更重要的作用是承担较高层次的信息处理和解释工作。具体来说，它要帮机器人应对以下几方面的精度需求：机器人本体运动学的精度、机器人外围设备的精度、工件本身的精度，以及作业空间里不断变化的环境条件。

按照传感器安装的位置和作用对象，可以分为\textbf{内部传感器}和\textbf{外部传感器}两类。内部传感器让\textbf{机器人感知自己的内部状态}，比如速度、位置、力和力矩；外部传感器则\textbf{负责获取机器人所处环境的信息}，包括感知被操作对象和周围环境的状态、识别和判断工件的性质、测量各种物理量，以及分析周围场景。

外部传感器按照工作原理又可以分为接触式和非接触式两大类。接触式传感器包括触觉类（如开关、触点矩阵、面式开关）和力/力矩类（用于测量夹持力、力和力矩）；非接触式传感器则包括视频光学类（二维图像处理、灰度处理、三维立体视觉）、感应/电容/磁性类（接近开关、测距仪、焊缝识别）、光学测距类（光电门、光切法、一维测距、二维/三维扫描仪），以及超声波类（接近开关、测距仪、扫描仪）。整个传感器分类结构如图 \ref{fig:sensorstruktur} 所示。

\begin{figure}[htbp]
  \centering
  \includegraphics[width=0.9\textwidth]{sensorstruktur.png}
  \caption{传感器分类结构}
  \label{fig:sensorstruktur}
\end{figure}

\subsubsection{接触式传感器}

触觉传感器通过直接接触物体来获取环境信息，最常见的例子是集成在夹爪里的二值开关，抓取动作完成后它会给出零件已成功抓取的信号，一旦抓取失败就触发重试或者错误处理流程；另外像安全防护用的开关垫也属于这一类。触觉传感器还有一种典型应用是焊枪嘴：把喷嘴做成可以小幅度摆动的结构，内部装两个微动开关，喷嘴沿焊枪轴线方向可以偏移最多6毫米，垂直方向最多25毫米，一旦喷嘴偏移，微动开关就被触发，从而在焊接开始前探测出焊缝的起点和终点，如果焊接过程中发生碰撞，还能立即中断焊接。这种触觉传感器的搜索速度大约是每分钟2米，搜索精度大约是正负0.2毫米。除了这种摆动式喷嘴，实际中还会用球形触头、圆盘形触头，以及专门用于内角焊缝的角部触头。触觉传感器的优点是成本低、结构简单，而且不怕焊接过程中的光学干扰；缺点则是容易磨损、测量精度有限，也容易被污染影响。

触觉传感器还有一类是面状分布的触觉阵列，利用不同的物理转换原理来生成矩阵式的触觉信息，讲义举了压电PVDF薄膜的例子：这是一种带有明显压电效应的热塑性塑料薄膜，厚度大约在10到100微米之间，薄膜两面镀铝，通过局部蚀刻的方式在薄膜上做出矩阵式的感应点，只有上下两面之间仍保持电极化关系的区域才会起作用。

图 \ref{fig:beruhrendesensor} 中所示的\textbf{力/力矩传感器}相当于给了机器人触觉，把这种传感器装进机器人后，可以实现\textbf{精密零件的装配定位、自动化装配作业，以及由机器人自主完成的质量检测}。安装位置通常在机器人手腕和末端执行器之间，或者直接装在末端执行器内部，也可以集成到工件夹具中。这类传感器内部一般带有数字信号处理器，负责采集和过滤应变信号、计算出对应的力和力矩数值，再通过通信接口把结果传出去。讲义还展示了一种典型结构：两层叠放的改造过的辐条轮结构，工作原理是先把受到的力和力矩转化成弹性辐条的微小长度变化，再通过应变片把这种变化转成电信号，最后借助一个所谓的"解耦矩阵"，从各个传感元件的信号中解算出力和力矩的各个分量，这本质上是求解一个和具体结构相关的线性方程组。这类传感器普遍存在的问题是信号里常常叠加着振荡成分，影响信号质量，虽然可以通过后处理和滤波来改善，但往往会牺牲一部分信号的动态响应速度。

\begin{figure}[htbp]
  \centering
  \includegraphics[width=0.4\textwidth]{beruhrendesensor.png}
  \caption{力/力矩传感器}
  \label{fig:beruhrendesensor}
\end{figure}

\subsubsection{非接触式传感器}

视频光学传感器的工作方式是先拍摄工件场景的图像，再通过计算机做图像处理，从而判断物体的种类、数量和位置。一套典型的图像处理系统通常需要先经历一个示教阶段：把要识别的物体以不同姿态呈现给系统，从中总结出效果较好的预处理、物体分割和特征计算算法，尽量用少而准确的特征来描述物体，图像处理软件在整个系统里是核心环节，针对不同任务通常会配有模块化的软件包。

二值图像处理系统对光照条件的变化非常敏感，很多情况下需要配合专门的光照控制装置；灰度图像处理系统的可靠性则明显更好。常见的照明方式包括透射光和反射光两种：透射光用来呈现透明物体的内部结构，或者拍出不透明物体的轮廓（需要有透明的衬底）；反射光则用来突出物体的表面结构，根据入射角的不同又可以分为漫反射光和定向反射光。漫反射光照明比较均匀，物体转动或者在承载面上平移时反射效果变化不大，因此常用于工件识别；定向反射光则能更好地呈现工件表面的纹理结构。为了消除运动模糊，有时还会用与生产节拍同步的闪光（多为红外光）。反射光还可以做成局部投影的形式，比如点状、线状或条纹状的图案，甚至是不规则或者面编码的图案投射到物体表面，这样就能实现三维形状的采集——投出的线条越细，说明测量距离越远，这种照明技术被称为光切法。

视频光学传感器另一个重要应用方向是物体识别，如图 \ref{fig:nichtberuhrendesensor} 所示。要让生产过程高效运转，一个基本前提是能在最短时间内准确无误地识别出物体，这就是自动识别（AutoID）技术要解决的问题。对一套自动识别系统的要求包括：在给定条件下保证足够高的读取可靠性、满足输送环节所需的读取速度、能够生成足够多样的识别标记、读取距离能够根据实际情况调整、与供应链上其他环节保持兼容，同时还要控制识别系统和相关设备的成本。识别所依据的特征既可以是物体本身固有的属性，比如颜色、重量、体积、材质、长宽高、包装材料、温度等，通过组合多个特征进行分类，进而完成物体的分拣；也可以是人工添加的识别标记，比如各种条码、二维码（QR码、矩阵码等），这些标识符可以存储在纸张、磁卡、打孔纸带或者RFID标签等不同载体上，再通过条码扫描器、二维相机等专门的读取设备来读取。另外还有一种识别技术是RFID（射频识别），它通过感应或者电磁波，在不需要直接光学接触的情况下就能非接触地采集和传输二进制编码的数据。

\begin{figure}[htbp]
  \centering
  \includegraphics[width=0.4\textwidth]{nichtberuhrendesensor.png}
  \caption{用于物体识别的视频光学传感器}
  \label{fig:nichtberuhrendesensor}
\end{figure}

感应式传感器大多利用涡流原理工作：一个开口的壳形磁芯是LC振荡电路的一部分，用来产生高频振荡。一旦有导电物体进入这个向外发散的磁场，物体内部就会产生涡流，这个涡流会从振荡回路里"抽走"一部分能量，导致振荡幅度下降，这个变化会被电子电路检测出来，再通过和一个固定阈值比较，就可以得到一个距离信号。常见的环形线圈感应器的工作范围大约是线圈半径的0.01到1倍。这种测量原理在机器人领域常用于距离极限开关、位移传感器和脉冲发生器。感应式传感器在实际应用中主要用于焊接机器人的焊缝跟踪，尤其适合对接焊缝和搭接焊缝这类形状规则的焊缝，因为结构简单、成本较低，常被用在焊缝起点检测、跟踪直线或轻微弯曲焊缝这类相对简单的任务上。但如果焊缝弯曲程度比较大，感应式传感器一般就不适用了（磁场较强时传感器需要提前探测，容易受影响），另外焊接飞溅和夹具也会产生歧义信号，进一步限制了它的使用场景。

\subsection{面向工业4.0数字化的3D测量技术}

数字化的本质是\textbf{把物体表面用离散化的方式重建出来}，\textbf{测量和检测数据能反映产品和加工过程的状态是质量保障的基础}，而数据所包含的信息量高低很大程度上取决于所用的传感系统。

激光类3D测量系统（全站仪、激光跟踪仪、激光扫描仪）在原理上都是\textbf{极坐标测量}：分别测出到目标点的距离和两个方向角，再合成为三维坐标。

距离测量最基本的三种方法是脉冲飞行时间法、相位比较法和三角测量法。
\begin{itemize}
    \item 脉冲飞行时间法通过测量激光脉冲往返的飞行时间来计算距离，是激光系统最主要的测距方式；
    \item 相位比较法通过观测和分析持续发射的载波信号的相位移完成测距，部分激光设备也采用这种方式，通常精度更高但量程受限；
    \item 三角测量法则是在固定基线的前提下，通过观测角度的变化来推算距离。
\end{itemize}
角度测量方面，水平角和垂直角由\textbf{角度编码器（Winkelencoder）}测得：既可以采用固定刻度盘的静态方法（编码轨迹法、增量法、插值法），也可以采用旋转刻度盘的动态方法（把方向测量转化为时间测量），此外还有不需要刻度盘的激光陀螺仪。

需要注意这两类系统在原理上的分工：激光系统依靠"测距 + 测角"直接得到点的极坐标，而相机系统和混合系统（摄影测量、条纹光投影）则依靠\textbf{三角测量法}获取深度信息——通过已知基线上的两个（或多个）视角对同一物点的交会来反算其空间坐标。

在极坐标式、逐点测量的系统中，比较典型的代表是全站仪和激光跟踪仪。全站仪的测程可以达到3500米，精度约为0.6毫米加上百万分之一的比例误差，主要用于对超大型结构做中等精度要求的测量；激光跟踪仪的测量体积可达120米，精度约为15微米再加上每米6微米的比例误差，适合对大型结构做非常高精度的测量。这类设备的典型应用包括船舶建造中的尺寸质量检验（比如在船体分段合拢过程中逐点检查节点位置、给出分段与理论位置的偏差以便进一步校正定位），以及风电叶片模型的几何质量控制（因为测量体积大而公差要求又很严格，对精度的要求非常高）。

面测量系统大多基于激光扫描仪，按测量体积和使用方式大致可以分为三类：手持式扫描仪测量体积较小，适合手动扫描自由曲面，采样速率可达每秒50万个点；近距离扫描仪测量体积中等，一般固定放置进行扫描，对固定场景的测量精度较高；地面激光扫描仪测量体积最大，以栅格方式扫描表面并同时记录反射光的强度，采样速率可以超过每秒100万个点，测量距离可以达到160米。地面激光扫描的典型应用包括船舶建造中对现状的数字化（用来指导后续的刮腻子等工序、提高预制加工的比例），以及对建筑毛坯结构的采集（作为设计幕墙构件、并根据实际轮廓做适配调整的依据）。

光学测量系统在结构上一般包含图像采集、图像传输、图像预处理和图像处理四个环节。图像采集依靠CCD或CMOS传感器完成矩阵式的离散采样和数字化；图像传输通过DVI、HDMI、WLAN等标准化信号输出；图像预处理运用各种数值方法改善图像质量、降低噪声、提高对比度；图像处理则负责自动提取特征或者做模式识别。其中一个重要的技术分支是摄影测量学，主要用于对静态场景做高精度测量：通过CCD矩阵完成图像测量，得到图像坐标；再通过"图像定向"，借助已知的定向点（比例尺）及其图像坐标，结合相机参数的近似值，反推出拍摄瞬间相机在空间物体坐标系中的位置和姿态；有了至少两张图像的图像坐标、至少四个定向点的物体坐标，再加上相机参数的近似值，就可以通过"前方交会"（多像三角测量）反算出物体上各点的坐标；最后还会通过"光束法平差"，把测量误差均匀分摊到所有参与计算的图像上，以提高整体精度。摄影测量学的典型应用包括吊车法兰的平面度检测（作为装配过程中的质量控制手段，同时也把摄影测量当作一种可移动的坐标测量设备来使用），以及船用夹层结构的变形分析（在多个加载阶段测量结构上按有限元网格划分的定点变形量，用来优化夹层结构的设计）。

除了摄影测量，另一种面测量方式是条纹光投影法：把某种图案投影到物体表面，再用立体相机系统拍摄成像，均匀分布的条纹图案在物体表面会发生形变，利用三角测量法计算出各个物体坐标点，再通过多边形化和渲染的方式把这些点重建成连续的曲面，还可以通过相移法等手段进一步提高精度。条纹光投影技术的典型应用包括对尺寸最大可达11米的螺旋桨做面状质量检测（采集并处理螺旋桨的几何形状，从不同预设视角拍摄后再拼接，把实测模型和理论CAD模型进行拟合比对，自动生成检测报告以保证质量的一致性），也可以用于批量生产开始阶段对自动化焊接车架的质量控制（拼接多次拍摄的结果，检查形状和位置公差的细节，并结合过程参数综合评估整个部件和加工精度）。

在实际应用中，单一测量系统的测量体积和采集速度往往有限，这时可以通过组合不同的传感器来弥补：让一个移动的传感器被另一个系统持续跟踪，作为参考的系统既可以是机械连接（比如测量臂），也可以是光学跟踪（比如激光跟踪仪、相机系统），这样组合起来的优点是可以实现外部参照，同时减少不同测量结果之间的拼接误差。

\subsection{机器人控制与传感器耦合应用}

传感器又是如何与机器人控制系统深度结合，实现自动化甚至自适应的加工的呢？

一种典型场景是自动化测量装置：条纹光投影系统安装在带直线轴的机器人上（总共可以有8个轴），可以在一个测量室内对大型工件（比如船用螺旋桨、钻孔工具）实现高精度的整体自动化测量；机器人的运动和末端执行器的位置也可以通过外部系统（比如激光跟踪仪）来测定和控制，测量系统会实时检查末端执行器的位置，并据此对机器人的运动进行修正。

另一种是基于传感器的机器人编程，这是一套用来自动生成加工程序的模块化系统：先在机器人作业空间内或者前置工位上，对工件做三维数据采集；再由控制计算机自动处理这些三维传感数据，从中判断出具体的加工任务；系统还可以选配图形化用户界面，把识别出的加工任务以图形方式展示出来，供操作人员做进一步的交互确认；最后通过后处理器把加工任务转换成设备专属的机器人程序，也可以借助宏指令来生成程序或者结合碰撞检测完成轨迹规划，生成的程序再传输给设备控制系统去执行。这种基于传感器的编程方式相比传统的手动编程有明显优势：能减少编程工作量，使用的是工件当下的实际数据和实际位置，不依赖预先的设计和工艺准备工作，程序质量稳定、不会受个人操作差异影响，能够在时间和空间上更好地融入生产流程，同时还能提升整体的柔性。

\paragraph{示例：自适应机器人编程}

自适应机器人编程任务是在直径最大约10米的螺旋桨上打定位孔，为后续的打磨工序做准备，要求孔的位置精度达到2毫米、孔深精度达到0.1毫米。如图 \ref{fig:adaptiveprogramming} 所示。这里的难点在于螺旋桨的具体摆放位置是未知的，同时精度要求又很高。解决办法是先用激光跟踪仪确定螺旋桨在机器人坐标系中的位置，再借助激光跟踪仪确定机器人上钻头相对于螺旋桨的位置，并据此修正每一个钻孔点的实际位置。在控制钻孔深度方面，钻孔的进给由一根独立的直线轴来实现，通过监测钻头电机电流的变化来判断是否已经开始切入材料，再结合三个激光测距传感器的信号，实时监控钻孔深度、并在达到设定深度时及时停止进给。

\begin{figure}[htbp]
  \centering
  \includegraphics[width=0.6\textwidth]{adaptiveprogramming.png}
  \caption{自适应机器人编程}
  \label{fig:adaptiveprogramming}
\end{figure}